%%%%%%%%%%%%%%%%%%%%%%%%%%%%%%%%%%%%%%%%%%%%%%%%%%%%%%%%%%%%%%%%%
%                                                               %
% Cette version repart de 03_Draft\bargaining_current-short.tex %
%                                                               %
%%%%%%%%%%%%%%%%%%%%%%%%%%%%%%%%%%%%%%%%%%%%%%%%%%%%%%%%%%%%%%%%%

% NB: Je supprime toutes les (grosses) parties qui étaient en commentaire dans la version précédente pour qu'on y voit plus clair.

Our findings on productivity dismiss the notion that workers are re-employed in ``bad'' (low productivity, low pay) firms. Instead, the reallocation pattern suggests that the wage cost of job displacement is driven by a reallocation toward low labor negotiating power firms, preventing workers from claiming a larger share of value-added. We explore this idea by studying collective bargaining in destination firms. 


We build two qualitative measures of firm-level labor relations, based on administrative data on collective agreements and workers' representatives elections. Since within-sector dynamics drive our results, the use of firm-level variables is key to investigate the loss in bargaining power associated with post-displacement reallocations.

We construct two outcomes: whether a firm signed at least one wage-related collective agreement (2005-2007) and firm-level average turnout at workers' representatives elections (2009-2012). \citet{carluccio2015trade} show that firms signing collective agreements are more likely to engage in rent-sharing. Higher election participation reflects better in-house labor relations. Both measures correlate positively with firm premia (Appendix Table~\ref{tab:firm_fe_nego} and Figure~\ref{fig:bscatter_firmfe_bargaining}).

We estimate our main model (equation \ref{eq:main_estimating_equation}) using these negotiation indices as outcomes. For collective agreements, we restrict the sample to individuals displaced from 2008 onward to measure agreements before layoffs (2005-2007). For election turnout, we use contemporary measurements (2009-2012) and show robustness to small layoffs where reverse causality is less likely. Appendix \ref{appsec:nego} discusses these aspects in detail.

The results, displayed in Figure \ref{fig:LE_accords}, show that destination firms are less likely to sign wage-related collective agreements. At $t^D+4$, re-employed displaced workers are 8 percentage points less likely to work in firms with wage agreements \textcolor{orange}{[5.9: Table 2, panel b, col. 4 at $t^D+4$ = $-0.059$]}. Appendix \ref{appsec:neg:results} indicates these firms are generally less likely to sign pay and working conditions agreements and more likely to conduct failed negotiations. The result is robust when controlling for firm size (using four legally-relevant thresholds).

Regarding professional elections, Figure \ref{fig:LE_elec_main} shows lower average turnout in destination firms. At $t^D+5$, the average turnout drops by 14 percentage points \textcolor{orange}{[8.9: Table 2, panel b, col. 5 at $t^D+5$ = $-0.089$]}. Appendix \ref{appsec:neg:results} shows destination firms face more failed elections, and the turnout loss is similar for small downsizing events, suggesting the effect is not solely due to the layoff itself.

(\textcolor{red}{pw: numbers to review}) \textcolor{orange}{[see the two orange notes above]}


